\section{MỆNH ĐỀ}

\subsection{TÓM TẮT LÝ THUYẾT}
\subsubsection{Mệnh đề, mệnh đề chứa biến}
\begin{itemize}
	\item [\iconMT] \indam{Mệnh đề:} Mệnh đề là một \indamm{ câu khẳng định} đúng hoặc sai. 
	\begin{boxkn}
		\begin{itemize}
			\item Câu khẳng định đúng gọi là mệnh đề đúng.
			\item Câu khẳng định sai gọi là mệnh đề \textbf{sai}.
			\item Một mệnh đề không thể vừa đúng, vừa sai.
			\item Những mệnh đề liên quan đến toán học được gọi là mệnh đề toán học.
		\end{itemize} 
	\end{boxkn}
	\item [\iconMT] \indam{Mệnh đề chứa biến:} Mệnh đề chứa biến là một câu khẳng định chứa biến nhận giá trị trong một tập $X$ nào đó và với mỗi giá trị của biến thuộc $X$ ta được một mệnh đề.
	
\end{itemize}
\subsubsection{Mệnh đề phủ định}
\begin{itemize}
	\item [\iconMT] Cho mệnh đề $P$. Mệnh đề \lq\lq không phải $P$ \rq\rq gọi là mệnh đề phủ định của $P$.
	\item [\iconMT] Chú ý: 
		\begin{boxkn}
	\begin{itemize}
		\item  Mệnh đề phủ định của $P$, kí hiệu là $\overline{P}$.
		\item  Nếu $P$ đúng thì $\overline{P}$ sai, nếu $P$ sai thì $\overline{P}$ đúng.
	\end{itemize}
	\end{boxkn}
\end{itemize}
\subsubsection{Mệnh đề kéo theo và mệnh đề đảo}
	Cho hai mệnh đề $P$ và $Q$.
		\begin{itemize}
			\item[\iconMT] \indam{Mệnh đề kéo theo:} Mệnh đề "Nếu $P$ thì $Q$" gọi là mệnh đề kéo theo, kí hiệu  $P\Rightarrow Q$.
				\begin{boxkn}
			\begin{itemize}
				\item  Mệnh đề $P\Rightarrow Q$ còn được phát biểu là "$P$ kéo theo $Q$" hoặc "Từ $P$ suy ra $Q$"
				\item  Mệnh đề này chỉ sai khi $P$ đúng và $Q$ sai.
			\end{itemize}
				\end{boxkn}
				\item[\iconMT] \indam{Mệnh đề đảo:} Cho mệnh đề $P\Rightarrow Q$. Khi đó, $Q\Rightarrow P$ gọi là mệnh đề đảo của $P\Rightarrow Q$.
		\end{itemize}
	\begin{khung4}{Lưu ý}
		Xét định lý dạng $P\Rightarrow Q$. Khi đó, ta có thể phát biểu định lý này theo một trong 2 cách sau:
		\begin{listEX}[2]
			\item[\ding{172}] $P$ là điều kiện đủ để có $Q$.
			\item[\ding{173}] $Q$ là điều kiện cần để có $P$.
		\end{listEX}
	\end{khung4}
\subsubsection{Mệnh đề tương đương}
\begin{itemize}
	\item [\iconMT] Cho hai mệnh đề $P$ và $Q$. Mệnh đề ``$P$ nếu và chỉ nếu $Q$''  gọi là  mệnh đề tương đương.
	\item [\iconMT] Chú ý:
	\begin{boxkn}
	\begin{itemize}
		\item  Mệnh đề ``$P$ nếu và chỉ nếu $Q$'' được kí hiệu là $P\Leftrightarrow Q$.
		\item 	Mệnh đề $P\Leftrightarrow Q$ đúng khi cả $P\Rightarrow Q$ và $Q\Rightarrow P$ cùng đúng.
	\end{itemize}
	\end{boxkn}
\end{itemize}
			\begin{khung4}{Lưu ý}
				Xét định lý dạng $P \Leftrightarrow Q$. Khi đó, ta có thể phát biểu định lý này theo một trong các cách sau:
				\begin{itemize}
					\item[\ding{172}] $P$ là điều cần và đủ để có $Q$.
					\item[\ding{173}] $P$ khi và chỉ khi $Q$.
					\item[\ding{174}] $P$ tương đương $Q$.
				\end{itemize}
			\end{khung4}
\subsubsection{Mệnh đề có chứa kí hiệu $\forall, \exists$}
	\begin{itemize}
		\item[\iconMT] Mệnh đề chứa kí hiệu với mọi: $\forall x \in X,\, P(x) $.
	\begin{boxkn}
		\begin{itemize}
			\item Mệnh đề này đúng khi tất cả các giá trị của $x \in X$ đều làm cho phát biểu $P(x)$ đúng.
			\item Nếu ta tìm được ít nhất một giá trị $x \in X$ làm cho $P(x)$ sai thì mệnh đề này \textbf{sai}.
		\end{itemize}
	\end{boxkn}
		\begin{vidu}
			Mệnh đề "Bình phương mọi số thực đều không âm" được viết là
			$\forall x \in \mathbb{R},\, x^2 \ge 0.$
		\end{vidu}
		\item[\iconMT] Mệnh đề chứa kí hiệu tồn tại: $\exists x \in X,\, P(x) $.
	\begin{boxkn}
		\begin{itemize}
			\item Mệnh đề này đúng khi ta tìm được ít nhất một giá trị của $x \in X$ làm cho phát biểu $P(x)$ đúng.
			\item Nếu tất cả giá trị của $x \in X$ đều làm cho $P(x)$ sai thì mệnh đề này \textbf{sai}.
		\end{itemize}
	\end{boxkn}
		\begin{vidu}
			Mệnh đề "Có một số tự nhiên mà bình phương của nó bằng 3" được viết là
			$\exists x \in \mathbb{N},\, x^2=3.$
		\end{vidu}
		\item[\iconMT] Phủ định của Mệnh đề chứa kí hiệu $\forall $, $\exists $.
	\begin{boxkn}
		\begin{itemize}
			\item Phủ định của mệnh đề ``$\forall x\in X,P\left( x \right)"$ là mệnh đề ``$\exists x\in X,\overline{P(x)}"$.
			\item Phủ định của mệnh đề ``$\exists x\in X,P\left( x \right)"$ là mệnh đề ``$\forall x\in X,\overline{P(x)}"$.
		\end{itemize}
	\end{boxkn}
	\end{itemize}

	
\subsection{RÈN LUYỆN KĨ NĂNG GIẢI TOÁN}
\begin{dang}{Mệnh đề, phủ định của mệnh đề}
	\begin{enumerate}[\iconMT]
		\item \indam{Mệnh đề:}
		\begin{listEX}[1]
			\item [\ding{172}] Khẳng định đúng là mệnh đề đúng, khẳng định sai là mệnh đề sai.
			\item [\ding{173}] Câu không phải là câu khẳng định hoặc câu khẳng định mà không có tính đúng-sai đều không phải là mệnh đề.
		\end{listEX}
		\item \indam{Mệnh đề phủ định:} Cho mệnh đề $P$. 
		\begin{listEX}[1]
			\item [\ding{172}] Mệnh đề phủ định của $P$, kí hiệu là $\overline{P}$.
			\item [\ding{173}] Nếu $P$ đúng thì $\overline{P}$ \textbf{sai}; $P$ \textbf{sai} thì $\overline{P}$ đúng.
		\end{listEX} 
	\end{enumerate}	
\end{dang}

\begin{vd}
	Các câu sau đây, câu nào là mệnh đề, câu nào không phải là mệnh đề? Nếu là mệnh đề hay cho biết mệnh đề đó đúng hay \textbf{sai}?
	\begin{listEX}[2]
		\item Không được đi lối này!
		\item Bây giờ là mấy giờ?
		\item $7$ không là số nguyên tố.
		\item $\sqrt{5}$ là số vô tỉ.
	\end{listEX}
	\loigiai{
		\begin{listEX}[1]
			\item Không phải là mệnh đề. Đây là câu cảm thán.
			\item Không phải mệnh đề. Đây là câu hỏi.
			\item Mệnh đề sai.
			\item Mệnh đề đúng
	\end{listEX}}
\end{vd}\dongcham{1}

\begin{vd}
	Xét tính đúng sai của các mệnh đề sau và lập mệnh đề phủ định của chúng.
	\begin{enumerate}
		\item[a)] $5$ là số nguyên tố.
		\item[b)] Phương trình $2x^2 - 3x + 1 = 0$ có nghiệm nguyên.
		\item[c)] Tổng ba góc trong một tam giác bằng $180^\circ$.
	\end{enumerate}
	\loigiai{
		\begin{enumerate}
			\item[a)] Đây là {mệnh đề đúng} vì $5$ chỉ có hai ước là $1$ và chính nó.\\
			\textbf{Mệnh đề phủ định:} ``$5$ không phải là số nguyên tố''.
			
			\item[b)] Giải phương trình $2x^2 - 3x + 1 = 0$, ta được các nghiệm là $x = 1$ và $x = \frac{1}{2}$. Vì phương trình có một nghiệm nguyên là $x = 1$ nên đây là {mệnh đề đúng}.\\
			\textbf{Mệnh đề phủ định:} ``Phương trình $2x^2 - 3x + 1 = 0$ không có nghiệm nguyên''.
			
			\item[c)] Đây là {mệnh đề đúng} (theo định lí tổng ba góc của một tam giác trong hình học phẳng).\\
			\textbf{Mệnh đề phủ định:} ``Tổng ba góc trong một tam giác không bằng $180^\circ$''.
		\end{enumerate}
	}
\end{vd}
\dongcham{1}
\begin{vd}%[Tex hóa SGK CD-CT,T7/22, TVN-001]%[0D1B1-3]
	Lập mệnh đề phủ định của mỗi mệnh đề sau và nhận xét tính đứng sai của mỗi mệnh đề phủ định đó.
	\begin{listEX}[1]
		\item $ A\colon $\lq\lq $ \dfrac{5}{1{,}2} $ là  một phân số\rq\rq;
		\item $ B\colon $\lq\lq Phương trình $ x^2+3x+2=0 $ có nghiệm\rq\rq;
		\item $ C\colon $\lq\lq $ 2^2+2^3=2^{2+3} $\rq\rq;
		\item $ D\colon $\lq\lq Số $ 2025 $ chia hết cho $ 15 $\rq\rq.
	\end{listEX}
	\loigiai{
		\begin{listEX}[1]
			\item Mệnh đề $ \overline{A}\colon $\lq\lq $ \dfrac{5}{1{,}2} $ không là một phân số\rq\rq. Mệnh đề $ \overline{A} $ đúng.
			\item Mệnh đề $ \overline{B}\colon $\lq\lq Phương trình $ x^2+3x+2=0 $ vô nghiệm\rq\rq. Mệnh đề $ \overline{B} $ sai.
			\item Mệnh đề $ \overline{C}\colon $\lq\lq $ 2^2+2^3\neq 2^{2+3} $\rq\rq. Mệnh đề $ \overline{C} $ đúng.
			\item Mệnh đề $ \overline{D}\colon $\lq\lq Số $ 2~025 $ không chia hết cho $ 15 $\rq\rq. Mệnh đề $ \overline{D} $ sai.
		\end{listEX}	
	}
\end{vd}
\dongcham{6}
\begin{dang}{Mệnh đề kéo theo, mệnh đề đảo}
\end{dang}

\begin{vd}
	Cho tam giác $ABC$. Xét hai mệnh đề 
	\begin{itemize}
		\item $P\colon$ \lq\lq Tam giác $ABC$ vuông\rq\rq;
		\item $Q\colon$\lq\lq Tam giác $ABC$ có $AB^2+AC^2=BC^2$ \rq\rq.
	\end{itemize}
Phát biểu các mệnh đề sau và cho biết mệnh đề sau đúng hay \textbf{sai}?
	\begin{listEX}[2]
		\item $P\Rightarrow Q$.
		\item $Q\Rightarrow P$.
	\end{listEX}
	\loigiai{
		\begin{listEX}[2]
			\item Mệnh đề $P\Rightarrow Q$ là ``Nếu tam giác $ABC$ vuông thì $AB^2+AC^2=BC^2$''.\\
			Mênh đề $P\Rightarrow Q$ sai vì chưa chắc tam giác đã vuông tại $A$.
			\item Mệnh đề $Q\Rightarrow P$ là ``Nếu tam giác $ABC$ có $AB^2+AC^2=BC^2$ thì tam giác vuông''.\\
			Mệnh đề $Q\Rightarrow P$ đúng (theo định lí Pitago).
	\end{listEX}}
\end{vd}\dongcham{3}

\begin{vd}
	Với mỗi cặp mệnh đề $P$ và $Q$ sau đây, hãy phát biểu mệnh đề $P \Rightarrow Q$ và xét tính đúng sai của nó.
	\begin{enumerate}[a)]
		\item $P$: ``$b^2 \ge 4ac$'';\\
		$Q$: ``Phương trình $ax^2 + bx + c = 0$ vô nghiệm'' ($a, b, c$ là ba số thực nào đó, $a \ne 0$).
		\item $P$: \lq\lq $a$ là số chia hết cho 6\rq\rq;\\
		$Q$: \lq\lq $a$ là số chia hết cho 3 \rq\rq.
		\item $P$: ``$a$ và $b$ là hai số chẵn'';\\
		$Q$: ``$a + b$ là số chẵn'' ($a, b$ là hai số tự nhiên).
		\item $P$: ``Tứ giác $ABCD$ có bốn cạnh bằng nhau'';\\
		$Q$: ``Tứ giác $ABCD$ là một hình vuông''.
	\end{enumerate}
	\loigiai{
		\begin{enumerate}
			\item[a)] Mệnh đề $P \Rightarrow Q$ phát biểu là: ``Nếu $b^2 \ge 4ac$ thì phương trình $ax^2 + bx + c = 0$ vô nghiệm'' ($a, b, c$ là ba số thực nào đó, $a \ne 0$).\\
			\textbf{Xét tính đúng sai:} Mệnh đề này là \textbf{sai}. \\
			\textit{Giải thích:} Xét phương trình bậc hai $ax^2 + bx + c = 0$ ($a \ne 0$) có biệt thức $\Delta = b^2 - 4ac$. Điều kiện $b^2 \ge 4ac$ tương đương với $\Delta \ge 0$. Khi đó, phương trình phải có nghiệm (nghiệm kép hoặc hai nghiệm phân biệt).
			
			\item[b)] Mệnh đề $P \Rightarrow Q$ phát biểu là: ``Nếu $a$ chia hết cho 6 thì $a$ chia hết cho 3''($a$ là một số tự nhiên)..\\
			\textbf{Xét tính đúng sai:} Mệnh đề này là \textbf{đúng}.\\
			\textit{Giải thích:} Vì $6 = 2 \cdot 3$, nên một số chia hết cho $6$ (tức là bội của $6$) thì chắc chắn sẽ chia hết cho $3$.
			
			\item[c)] Mệnh đề $P \Rightarrow Q$ phát biểu là: ``Nếu $a$ và $b$ là hai số chẵn thì $a + b$ là số chẵn'' ($a, b$ là hai số tự nhiên).\\
			\textbf{Xét tính đúng sai:} Mệnh đề này là \textbf{đúng}. \\
			\textit{Giải thích:} Tổng của hai số chẵn luôn là một số chẵn.
			
			\item[d)] Mệnh đề $P \Rightarrow Q$ phát biểu là: ``Nếu tứ giác $ABCD$ có bốn cạnh bằng nhau thì tứ giác $ABCD$ là một hình vuông''.\\
			\textbf{Xét tính đúng sai:} Mệnh đề này là \textbf{sai}. \\
			\textit{Giải thích:} Tứ giác có bốn cạnh bằng nhau là hình thoi. Hình thoi cần thêm điều kiện có một góc vuông (hoặc hai đường chéo bằng nhau) mới trở thành hình vuông.
		\end{enumerate}
	}
\end{vd}
\dongcham{7}
\begin{vd}
	Phát biểu mệnh đề đảo của các mệnh đề sau và xét tính đúng sai của mệnh đề đảo đó.
	\begin{enumerate}
		\item Nếu tam giác $ABC$ có $AB = AC$ thì tam giác $ABC$ cân.
		\item Nếu tam giác $ABC$ có $2$ góc bằng $60^\circ$ thì tam giác $ABC$ đều.
	\end{enumerate}
	\loigiai{
		\begin{enumerate}
			\item[a)] \textbf{Mệnh đề đảo:} ``Nếu tam giác $ABC$ cân thì tam giác $ABC$ có $AB = AC$''.\\
			\textbf{Xét tính đúng sai:} Mệnh đề đảo này là \textbf{sai}.\\
			\textit{Giải thích:} Tam giác $ABC$ cân thì có thể cân tại đỉnh $B$ (khi đó $BA = BC$) hoặc cân tại đỉnh $C$ (khi đó $CA = CB$), không nhất thiết phải có $AB = AC$ (chỉ xảy ra khi tam giác cân tại đỉnh $A$).
			
			\item[b)] \textbf{Mệnh đề đảo:} ``Nếu tam giác $ABC$ đều thì tam giác $ABC$ có $2$ góc bằng $60^\circ$''.\\
			\textbf{Xét tính đúng sai:} Mệnh đề đảo này là \textbf{đúng}.\\
			\textit{Giải thích:} Tam giác đều thì cả ba góc của nó đều bằng $60^\circ$, do đó hiển nhiên thỏa mãn việc có $2$ góc bằng $60^\circ$.
		\end{enumerate}
	}
\end{vd}
\dongcham{4}
\begin{dang}{Mệnh đề chứa kí hiệu với mọi, tồn tại}
	\begin{enumerate}[\iconMT]
		\item \indam{Tính đúng sai:}
		\begin{listEX}[1]
			\item [\ding{172}]  Mệnh đề chứa kí hiệu với mọi: $\forall x \in X,\, P(x) $.
				\begin{itemize}
					\item Mệnh đề này đúng khi tất cả các giá trị của $x \in X$ đều làm cho phát biểu $P(x)$ đúng.
					\item Nếu ta tìm được ít nhất một giá trị $x \in X$ làm cho $P(x)$ sai thì mệnh đề này \textbf{sai}.
				\end{itemize}
			\item [\ding{173}] Mệnh đề chứa kí hiệu tồn tại: $\exists x \in X,\, P(x) $.
					\begin{itemize}
					\item Mệnh đề này đúng khi ta tìm được ít nhất một giá trị của $x \in X$ làm cho phát biểu $P(x)$ đúng.
					\item Nếu tất cả giá trị của $x \in X$ đều làm cho $P(x)$ sai thì mệnh đề này \textbf{sai}.
				\end{itemize}
		\end{listEX}
		\item \indam{Phủ định của mệnh đề có dấu $\forall,\exists $:}
		\begin{listEX}[1]
			\item [\ding{172}] $\forall x\in X,P(x)$ thành $\exists x\in X,\overline{P(x)}$.
			\item [\ding{173}] $\exists x\in X,P(x)$ thành $\forall x\in X,\overline{P(x)}$.
		\end{listEX}
	\indam{Chú ý:} \indamm{Khi lấy phủ định, ta chú ý các vấn đề đối lập sau:}
	\begin{listEX}[1]
		\item [\ding{172}] Quan hệ $=$ thành quan hệ $\ne $, và ngược lại.
		\item [\ding{173}] Quan hệ $>$ thành quan hệ $\le $, và ngược lại.
		\item [\ding{174}] Quan hệ $\ge $ thành quan hệ $<$, và ngược lại.
		\item [\ding{175}] Liên kết "và" thành liên kết "hoặc", và ngược lại.
	\end{listEX}
	\end{enumerate}
\end{dang}

\begin{vd}%[0D1Y1-5]
	Sử dụng kí hiệu \lq\lq$\forall$\rq\rq \,để viết mỗi mệnh đề sau và xét xem mệnh đề đó là đúng hay sai?
	\begin{enumerate}
		\item $P\colon$\lq\lq Với mọi số thực $x, x^2+1>0$\rq\rq.
		\item $Q\colon$\lq\lq Với mọi số tự nhiên $n, n^2+n$ chia hết cho $6$\rq\rq.
	\end{enumerate}
	\loigiai{
		\begin{enumerate}
			\item $P\colon$\lq\lq Với mọi số thực $x, x^2+1>0$\rq\rq.\\
			Mệnh đề được viết là $P \colon \lq\lq\forall x \in \mathbb{R}, x^2+1>0$\rq\rq.\\
			Xét một số thực $x$ tùy ý, ta phải chứng tỏ rằng $x^2+1>0$.\\
			Thật vậy, ta có $x^2+1 \geq 1>0$.\\
			Vậy mệnh đề $P$ là mệnh đề đúng.
			\item $Q\colon$\lq\lq Với mọi số tự nhiên $n, n^2+n$ chia hết cho $6$\rq\rq.\\
			Mệnh đề được viết là $Q\colon\lq\lq \forall n \in \mathbb{N},\left(n^2+n\right) \,\vdots\, 6$\rq\rq.\\
			Để chứng minh mệnh đề $Q$ là sai, ta cần chỉ ra một giá trị cụ thể của $n$ để nhận được mệnh đề sai.\\
			Thật vậy, chọn $n=1$, ta thấy $n^2+n=2$ không chia hết cho $6$.\\
			Vậy mệnh đề $Q$ là mệnh đề sai.
		\end{enumerate}
	}
\end{vd}\dongcham{4}

\begin{vd}%[0D1Y1-5]
	Sử dụng kí hiệu \lq\lq$\exists$\rq\rq\, để viết mỗi mệnh đề sau và xét xem mệnh đề đó là đúng hay sai?
	\begin{enumerate}
		\item $M\colon$\lq\lq Có ít nhất một số thực $x$ sao cho $x^3=-8$\rq\rq.
		\item $N\colon$\lq\lq Tồn tại số nguyên $x$ sao cho $2x+1=0$\rq\rq.
	\end{enumerate}
	\loigiai{
		\begin{enumerate}
			\item $M\colon$\lq\lq Có ít nhất một  số thực $x$ sao cho $x^3=-8$\rq\rq.\\
			Mệnh đề được viết là $M\colon\lq\lq\exists x \in \mathbb{R}, x^3=-8$\rq\rq.
			Để chứng tỏ mệnh đề $M$ là đúng, ta cần chỉ ra một giá trị cụ thể của $x$ để nhận được mệnh đề đúng.\\
			Thật vậy, chọn $x=-2$, ta thấy $(-2)^3=-8$.\\
			Vậy mệnh đề $M$ là mệnh đề đúng.\\
			Mệnh đề $N\colon\lq\lq\exists x \in \mathbb{Z}, 2x+1=0$\rq\rq.
			\item $N\colon$\lq\lq Tồn tại số nguyên $x$ sao cho $2x+1=0$\rq\rq.\\
			Để chứng minh mệnh đề $N$ là sai, ta phải chứng tỏ rằng với số nguyên $x$ tùy ý thì $2x+1 \neq 0$.\\
			Thật vậy, xét một số nguyên $x$ tùy ý, ta có $2x+1 \neq 0$.\\
			Vì thế mệnh đề $N$ là mệnh đề sai.
		\end{enumerate}
	}
\end{vd}\dongcham{4}

\begin{vd}
	Xét tính đúng sai của mệnh đề sau và nêu mệnh đề phủ định của nó.
	\begin{tasks}(2)
		\task $\exists x\in \mathbb{Z},x^2=3$.
		\task $\forall n\in {\mathbb{N}}^{*}:2^n+3$ là một số nguyên tố.
		\task $\forall x\in \mathbb{R},x^2+4x+5>0$.
		\task $\forall x\in \mathbb{R},x^4-x^2+2x+2\ge 0$.
	\end{tasks}
	\loigiai{
		\begin{listEX}
			\item Ta có $x^2=3\Leftrightarrow x=\pm \sqrt{3}$. Vì $\pm \sqrt{3}\notin \mathbb{Z}$ nên mệnh đề đã cho sai.\\
			Mệnh đề phủ định là $\overline{P(x)}:``\forall x\in \mathbb{Z},x^2\ne 3"$.
			\item Với $n= 5$ thì $2^n+3=2^5+3=35$, số này chia hết cho $5$ (không nguyên tố). Do đó mệnh đề đã cho sai.\\
			Mệnh đề phủ định là $\overline{P(n)}:``\exists n\in {\mathbb{N}}^{*}:2^n+3$ không là một số nguyên tố''.
			\item Mệnh đề đúng vì $x^2+4x+5={\left(x+2\right)}^2+1>0, \forall x\in \mathbb{R}$.\\
			Mệnh đề phủ định là $\overline{P(x)}:``\exists x\in \mathbb{R},x^2+4x+5\le 0"$.
			\item Do $x^4-x^2+2x+2={\left(x^2-1\right)}^2+{\left(x+1\right)}^2\ge 0, \forall x\in \mathbb{R}$ nên mệnh đề đã cho đúng.\\
			Mệnh đề phủ định là $\overline{P(x)}:``\exists x\in \mathbb{R},x^4-x^2+2x+2<0"$.
		\end{listEX}
	}
\end{vd}\dongcham{7}

\begin{vd}
	Nêu mệnh đề phủ định của các mệnh đề sau và cho biết tính đúng sai của mệnh đề phủ định đó.
	\begin{listEX}[2]
		\item $A:``\exists n\in \mathbb{N},n^2+3$ chia hết cho $4$''.
		\item $B:``\exists x\in \mathbb{N}$, $x$ chia hết cho $x+1$''.
	\end{listEX}
	\loigiai{
		\begin{listEX}
			\item $\overline{A}:``\forall x\in \mathbb{N},n^2+3$ không chia hết cho $4$''. Mệnh đề này sai.
			\item $\overline{B}:``\forall x\in \mathbb{N},x$ không chia hết $x+1$''. Mệnh đề này sai.
		\end{listEX}
	}
\end{vd}\dongcham{7}

\boxmini{BÀI TẬP TỔNG HỢP TỰ LUYỆN}

\begin{baitap}%[0D1Y1-2]
	Trong các mệnh đề toán học sau đây, mệnh đề nào đúng? Mệnh đề nào sai?
	\begin{itemize}
		\item [a)]$P\colon$\lq\lq Tổng hai góc đối của một tứ giác nội tiếp bằng $180^{\circ}$\rq\rq.
		\item [b)]$Q\colon$\lq\lq $7$ là số chính phương\rq\rq.
		\item [c)]$R\colon$\lq\lq $1$ là số nguyên tố\rq\rq.
	\end{itemize}
	\loigiai
	{
		Mệnh đề $P$ là mệnh đề đúng. \\
		Mệnh đề $Q$ và $R$ là mệnh đề sai. \\
	}
\end{baitap}\cham{9}

\begin{baitap}%[0D1Y1-2]
	Xét tính đúng sai của mỗi mệnh đề sau
	\begin{tasks}(2)
		\task $1993$ chia hết cho $3$.
		\task $\sqrt{12}$ là một số hữu tỉ.
		\task $9$ là một số chính phương.
		\task $|-1997|\leqslant0$.
		\task $\pi<\dfrac{10}{3}$. 
		\task Phương trình $3x+7=0$ có nghiệm.
	\end{tasks}
	\loigiai
	{
		\begin{enumerate}[a)]
			\item Mệnh đề \lq\lq $1993$ chia hết cho $3$\rq\rq\ là mệnh đề sai vì $1993$ chia $3$ dư $1$.
			\item Mệnh đề \lq\lq $\sqrt{12}$ là một số hữu tỉ\rq\rq\ là mệnh đề sai vì $\sqrt{12}$ là một số vô tỉ.
			\item Mệnh đề \lq\lq $9$ là một số chính phương\rq\rq\ là mệnh đề đúng vì $\sqrt{9}=3$.
			\item Mệnh đề \lq\lq $|-1997|\leqslant0$\rq\rq\ là mệnh đề sai vì $|-1997|=1997>0$.
			\item Mệnh đề\lq\lq $\pi<\dfrac{10}{3}$\rq\rq\ là mệnh đề đúng.
			\item Mệnh đề \lq\lq Phương trình $3x+7=0$ có nghiệm\rq\rq\ là mệnh đề đúng vì $3x+7=0 \Leftrightarrow x=-\dfrac{7}{3}$.
		\end{enumerate}
	}
\end{baitap}\cham{9}

\begin{baitap}%[0D1K1-4]
	Sử dụng thuật ngữ ``điều kiện đủ'' để phát biểu các định lí sau.
	\begin{enumerate}
		\item Nếu $a$ và $b$ là hai số hữu tỉ thì tổng $a+b$ là số hữu tỉ.
		\item Nếu hai tam giác bằng nhau thì chúng có diện tích bằng nhau.
		\item Nếu một số tự nhiên có chữ số tận cùng là chữ số 5 thì nó chia hết cho $5$.
	\end{enumerate}
	\loigiai{
		\begin{enumerate}
			\item Điều kiện đủ để tổng $a+b$ là số hữu tỉ là cả hai số $a$ và $b$ đều là số hữu tỉ.\\
			Hoặc: $a$ và $b$ là hai số hữu tỉ là điều kiện đủ để tổng $a+b$ là số hữu tỉ.
			\item Điều kiện đủ để hai tam giác có diện tích bằng nhau là chúng bằng nhau.\\
			Hoặc: hai tam giác bằng nhau là điều kiện đủ để chúng có diện tích bằng nhau.
			\item Điều kiện đủ để một số chia hết cho 5 là số đó tận cùng bằng 5.\\
			Hoặc: một số tự nhiên có chữ số tận cùng là chữ số 5 là điều kiện đủ để nó chia hết cho 5.
			
		\end{enumerate}
	}
\end{baitap}\cham{15}

\begin{baitap}%[0D1B1-5]
	Dùng kí hiệu $\forall$ hoặc $\exists$ để viết các mệnh đề sau và xét tính đúng sai của chúng.
	\begin{tasks}(1)
		\task Mọi số thực khác $0$ nhân với nghịch đảo của nó bằng $1$.
		\task Có số tự nhiên mà bình phương của nó bằng $20$.
		\task Bình phương của mọi số thực đều dương.
		\task Có ba số tự nhiên khác $0$ sao cho tổng bình phương của hai số bằng bình phương của số còn lại.
	\end{tasks}
	\loigiai{
		\begin{enumerate}[a)]
			\item $\forall x \in \mathbb{R},x \ne 0,x \cdot \dfrac{1}{x}=1$. Mệnh đề đúng.
			\item  $\exists x \in \mathbb{N},x^2=20$. Mệnh đề \textbf{sai}.
			\item  $\forall x \in \mathbb{R},x^2>0$. Mệnh đề \textbf{sai}.
			\item  $\exists x,y,z \in \mathbb{N}^*,x^2+y^2=z^2$. Mệnh đề đúng.
	\end{enumerate}}
\end{baitap}\cham{17}

\begin{baitap}%[Nguyễn Cường- BG Toán 10]%[0D1B1-5]
	Lập mệnh đề phủ định của mỗi mệnh đề sau
	\begin{tasks}(2)
		\task $\forall x \in \mathbb{R},|x| \geq x$.
		\task $\exists x \in \mathbb{R}, x^2+1=0$.
	\end{tasks}
	\loigiai{
		\begin{enumerate}[a)]
			\item Phủ định của mệnh đề \lq\lq$\forall x \in \mathbb{R},|x| \geq x$\rq\rq\,là mệnh đề \lq\lq$\exists x \in \mathbb{R},|x|<x$\rq\rq.
			\item Phủ định của mệnh đề \lq\lq$\exists x \in \mathbb{R}, x^2+1=0$\rq\rq\,là mệnh đề \lq\lq$\forall x \in \mathbb{R}, x^2+1 \neq 0$\rq\rq.
		\end{enumerate}
	}
\end{baitap}\cham{10}


\begin{baitap}%[Nguyễn Cường- BG Toán 10]%[0D1B1-5]
	Lập mệnh đề phủ định của mỗi mệnh đề sau và xét tính đúng sai của mỗi mệnh đề phủ định đó
	\begin{tasks}(2)
		\task $\forall x \in \mathbb{R}, x^2 \neq 2x-2$.
		\task $\forall x \in \mathbb{R}, x^2 \leq 2x-1$.
		\task $\exists x \in \mathbb{R}, x+\dfrac{1}{x} \geq 2$.
		\task $\exists x \in \mathbb{R}, x^2-x+1<0$.
	\end{tasks}
	\loigiai{
		\begin{enumerate}[a)]
			\item $\exists x \in \mathbb{R}, x^2=2x-2$.\\
			Mệnh đề này sai vì phương trình $x^2-2x+2=0$ vô nghiệm trên tập số thực.
			\item $\exists x \in \mathbb{R}, x^2 > 2x-1$.\\
			Mệnh đề này đúng vì với $x=2$ thì $2^2>2\cdot 2-1$.
			\item $\forall x \in \mathbb{R}, x+\dfrac{1}{x}<2$.\\
			Mệnh đề này sai vì với $x=1$ thì $1+\dfrac{1}{1}=2$.
			\item $\forall x \in \mathbb{R}, x^2-x+1\ge 0$.\\
			Mệnh đề này đúng vì $x^2-x+1=\left(x-\dfrac{1}{2}\right)^2+\dfrac{3}{4}> 0$ với mọi $x\in\mathbb{R}$.
		\end{enumerate}	
	}
\end{baitap}\cham{21}
